% Response to the Associate Editor -- standalone letter.
% Compiles independently of the manuscript; shares no files with paper_dtrap.tex.
% In Overleaf, set Menu -> Main document to choose which of the two to build.
\documentclass[11pt,a4paper]{article}

\usepackage[margin=1in]{geometry}
\usepackage[T1]{fontenc}
\usepackage{lmodern}
\usepackage{microtype}
\usepackage{parskip}
\usepackage{booktabs}
\usepackage{longtable}
\usepackage{enumitem}
\usepackage{xcolor}
\usepackage[hidelinks]{hyperref}

% Editor quotations, set apart from our reply.
\definecolor{quoterule}{gray}{0.55}
\newenvironment{editorquote}
  {\par\medskip\noindent\begin{list}{}{%
     \setlength{\leftmargin}{2em}\setlength{\rightmargin}{1em}}\item[]%
     \color{black!72}\itshape}
  {\end{list}\medskip}

\newcommand{\pc}{$p_c$}
\setlength{\emergencystretch}{2em}

\title{\vspace{-2.2em}\textbf{Response to the Associate Editor}}
\date{}
\author{}

\begin{document}
\maketitle
\vspace{-3.5em}

\noindent
\begin{tabular}{@{}ll@{}}
\textbf{Manuscript ID} & COL-26-0021 \\
\textbf{Title} & Governing Societies of AI Agents: Committed Minorities, Model
                 Diversity, \\ & and Scale as Levers of Multi-Agent AI Governance \\
\textbf{Authors} & Gopal Kalyanaraman, Vijay K.\ Madisetti \\
\textbf{Journal} & \emph{Collective Intelligence} \\
\end{tabular}

\bigskip\hrule\bigskip

Dear Dr.\ Becker,

Thank you for the decision on COL-26-0021 and, in particular, for the unusual step
of returning the manuscript with a specific structural diagnosis rather than
sending it to referees in a form they could not assess. That was the right call,
and acting on it improved the paper considerably more than we anticipated---for
reasons we set out in Part~B below.

We respond first to your comment point by point. Part~B then reports the additional
experimental work the revision carries, which goes beyond what you asked for and
which we think substantially strengthens the paper's contribution.

\section*{Part A --- Response to the Associate Editor's comment}

\begin{editorquote}
``your paper is missing sufficiently detailed methods description, e.g.\ a methods
section\ldots{} I believe both that this paper holds promise, and also that
reviewers will be unable to assess the paper in its present form.''
\end{editorquote}

We agree, and we accept the diagnosis without reservation. The original submission
described \emph{what} each testbed measured but not \emph{how}, which left a referee
unable to judge whether the measurements supported the claims. We have added a
\textbf{Methods} section and a detailed \textbf{Appendix~B}. Following your
suggestion below, the body carries a compact methods overview (${\sim}820$ words)
and the full protocols, prompt templates and parameter table live in the appendix,
so the body stays readable while the paper remains conceptually replicable on its
own.

\begin{editorquote}
``Please note that making your reproducibility code available to reviewers would
not suffice for methods. The paper must stand on its own, and allow for full
conceptual replication, without reference to any supplied reproducibility code.''
\end{editorquote}

Understood, and the new material is written to that standard. Appendix~B reproduces
\textbf{every prompt template verbatim}, together with the neutral-label pool, the
parsing rules, and a parameter table giving population size, rounds or horizon,
seed count and committed-fraction grid for every study. A reader can reconstruct
each experiment from the paper alone.

\begin{editorquote}
``how you constructed the LLM interaction''
\end{editorquote}

\emph{Methods $\to$ Agents, Models, and Providers} defines an agent as a stateless
chat-completion call carrying its current state and locally visible context, with
all state held by the simulator rather than the model. This is what makes
population \emph{composition} a free variable and is the mechanism underlying every
homogeneous-versus-heterogeneous comparison in the paper.

\begin{editorquote}
``how you prompted the LLMs''
\end{editorquote}

Each paradigm's subsection states its prompting scheme, and Appendix~B gives the
templates verbatim. We draw particular attention to \emph{Methods $\to$ Binary
Agreement}, which documents a control absent from the original submission: the
literal tokens ``A'' and ``B'' carry a large, model-specific prior that dominates
the social dynamics, so opinions are realised as randomly drawn neutral nouns and
every configuration is run in both directions. We report the residual bias
statistic ($h \approx 0$ in all conditions).

\begin{editorquote}
``how you measured outcomes to get the reported results''
\end{editorquote}

Each subsection now states its estimator explicitly---the order parameter and
logistic fit with bootstrap intervals for the capture threshold; coverage, error
correlation and realised accuracy for the ensemble studies; survival, stock
trajectory and yield for the commons.

\emph{Methods $\to$ Data-Quality Auditing} adds a methodological control we consider
important enough to generalise beyond this paper. Multi-agent simulations degrade
\emph{silently} under partial API failure, because the natural fallback for a failed
decision is itself a plausible observation: in the binary-agreement model an
unchanged opinion is indistinguishable from an agent that deliberately held its
ground, so a dead endpoint renders as a population that perfectly resisted capture;
in the commons the fallback is the sustainable harvest, so a dead endpoint renders
as textbook cooperation. We now treat failed and unparseable responses as data loss
rather than observations, record the fallback rate for every run, and discard runs
above a small threshold. This was not precautionary: \textbf{three of six re-run
naming-game conditions proved silently corrupt}, and each had produced a plausible,
publishable-looking number. \emph{Threats to Validity} quantifies the bias---a
10\%-degraded run gave \pc{} $=0.0795$ against a clean $0.0985$, a 19\% error in the
direction that resembles a genuine finding. We have since applied the same
discipline to our own new experiments: every run reported in this revision records
its API-failure and parse-failure rates, and any arm above 5\% loss is marked
unreliable and discarded rather than reported.

\begin{editorquote}
``You may wish to include a general methods overview in the main text of the paper,
with more detailed methods in an Appendix.''
\end{editorquote}

Adopted exactly as suggested.

\begin{editorquote}
``I also strongly encourage you to make replication code available to reviewers.''
\end{editorquote}

All code, all result files, and the figure-generation scripts are public at
\url{https://github.com/kalyanar/governing-societies-of-ai-agents}. Every number in
the manuscript is traceable to a JSON file in \texttt{code/results/}. The repository
README documents setup, the API credentials each provider requires, and---for each
experiment---the claim it was built to test. Result files from runs that failed the
data-quality audit are retained rather than deleted, so that our exclusions are
auditable.

\section*{Part B --- Additional experiments completed for this revision}

Writing the Methods section to the standard you set required us to specify each
procedure precisely enough to re-run it. Having done so, we re-ran them, and then
extended them. This revision reports \textbf{ten new experiments}, and the paper is
materially stronger for it: the central result is measured at higher resolution, a
claim we had been prepared to withdraw turns out to hold under a fair test, the
cost result is re-established under considerably stricter conditions, and four
findings are new. Three previously reported numbers are superseded by better
measurement; we identify each explicitly at the end of this part.

\paragraph{1. The capture threshold generalises across content (new).}
Our original claim that the threshold carries from opinion formation to cooperative
behaviour required a fair test, which the commons study turned out not to provide
(item~2). We therefore held the coordination mechanism fixed---identical update
rule, population, grid, estimator and counterbalancing, through the same code
path---and varied only what is being agreed: neutral codewords; two conservation
rules of equal merit; and a restrained versus a greedy harvest norm. At $N=48$ with
five seeds and both push directions (70 episodes per condition, 0\% parse loss),
\textbf{the tipping point does not move}: \pc{} $=0.128$, $0.136$, $0.131$, spread
$0.0087$, every bootstrap interval overlapping. The asymmetric condition shows
residual directional bias $h=0.000$, so the identical threshold is not an artifact
of the two options being interchangeable in substance. The threshold is a property
of how agreement propagates, not of its subject matter, and it therefore transfers
to cooperation organised as a convention with rival stable states---the structure of
most norms, standards and protocols an agent society would negotiate.

\paragraph{2. Why a commons behaves differently, and a law that recovers our
${\sim}25\%$ (new).}
The fishery updates as $S \to \min(2(S-h),K)$: harvesting at or below $S/2$ pins the
stock at capacity and anything above sends it down geometrically, with no
intermediate rest point. The system has one attractor and an absorbing death state.
A tipping point requires \textbf{bistability}---two states each stable under
perturbation, separated by a barrier---so a depletion process cannot exhibit one at
any adversary share; it exhibits a \emph{rate}. Inverting the collapse-time law
gives the tolerance an observer reports after watching for $H$ months,
\[
  p_c(H) \;=\; 1-\left(S_c/S_0\right)^{1/H},
\]
which returns \textbf{25.9\% at $H=10$}---reproducing our published figure to within
a percentage point. Our ${\sim}25\%$ was therefore correct for the horizon at which
it was measured; what was missing was the horizon itself. We now report the law
rather than the point estimate, together with the general observation that any
capture threshold quoted without its observation window overstates robustness by a
computable amount. We note that the human result usually cited alongside ours
(Centola et al.\ 2018, ${\sim}25\%$) concerns social \emph{conventions}, which are
bistable, and is thus structurally a naming game rather than a commons.

\paragraph{3. The cost--quality result, re-established under stricter conditions
(new design).}
Our original comparison relied on a price table and an unconstrained judge. We have
replaced it with a head-to-head in which \textbf{both arms use the same judge model,
the same token budget and the same items}, so the comparison cannot be confounded by
model identity or by handicapping either side. A panel of three zero-cost models
reasons; the judge may then only return an index into their candidates. Result:
\textbf{$0.938$ at \$0.454 against $0.912$ at \$1.148 for that same model reasoning
alone---$2.53\times$ cheaper}; a panel of 24--31B models gives $1.88\times$. We
report this as \textbf{parity at a $2.5\times$ discount rather than as an
improvement}, since the accuracy margin is not significant (McNemar exact,
$p=0.625$), and we say so in the text. The saving is a token accounting rather than
a sampling estimate, and because both arms use one model the ratio is invariant to
that model's price---which removes the dependence on a price table that a referee
would otherwise have to date-check.

\paragraph{4. Two preconditions that make the architecture deployable (new).}
The substitution in item~3 is not universally available, and we can now say exactly
when it is. The panel's coverage must exceed the arbiter's solo accuracy: across
four panels spanning a factor of fifteen in member capability, this rule decides
every case, including the two where the architecture loses. And the panel must be
\textbf{competence-matched}---with members at $0.912/0.650/0.637$ the selector simply
tracks the leader and realises $\hat v = 0.00$, so diversity of lineage without
comparable competence buys nothing. Relatedly, blind majority voting fell
\emph{below} the panel's own best member in three of the four panels: aggregation by
counting is not a weak form of arbitration but an actively harmful one wherever
competence is uneven. Both conditions are measurable before any system is deployed.

\paragraph{5. A decisive control for constrained arbitration (new).}
On items where \textbf{no} panel member is correct, a genuine selector must be wrong
every time; any correct answer there proves the mechanism is generating rather than
selecting. Pooling four independent arms gives 112 such items. \textbf{A constrained
selector was correct on 0 of 112; a free-form judge was correct on 26.} Under the
constrained instrument the competence-band prediction holds end to end, with a
competent selector realising $\hat v = 0.33$--$0.91$ of the available headroom across
four panels. This control also explains an artifact in our earlier verifier
numbers, discussed below.

\paragraph{6. Verification is easier than generation, and the gap closes with
capability (new).}
Every verifier we tested selects better than it solves---eleven models from 2B to
frontier, ratios from $1.95\times$ down to $1.03\times$, declining monotonically with
capability. This is the mechanism that makes an arbitration architecture possible at
all, and it carries a caution we now state explicitly: because \emph{every} model
selects better than it solves, a candidate verifier beating its own solo score is no
evidence of fitness. The relevant comparison is against the panel's best member, on
which every sub-frontier verifier we tried failed.

\paragraph{7. The capability floor is a property of the role (new).}
Sweeping eleven verifiers from 2B to frontier against a fixed panel, every model up
to and including a current cheap frontier model \emph{selected worse than the
panel's own best member}. A model that clears the coordination floor at 8B is
nowhere near the arbitration floor. A capability gate is therefore a (role, floor)
pair, and a governance regime that admits agents on a single threshold will
misallocate them across roles even when that threshold is measured rather than
assumed.

\paragraph{8. The central result at higher resolution (strengthened).}
Our original sweep used $N=16$ on a grid whose only interior point was $p=0.1$, so
the entire transition fell between two grid points. Re-running all five lineages at
$N=64$, with grid points on exact multiples of $1/N$ and seven resolving points
inside the transition, \textbf{strengthens the paper's central claim}: the lineages
pool to \pc{} $=0.0959 \pm 0.0019$, with the analytic $0.0979$ inside every
individual interval and a cross-lineage spread seven times tighter.

\paragraph{Measurements updated by the above.}
Three figures from the original submission are replaced by the work described here.
\emph{(i)} The five per-lineage thresholds were estimated on an $N=16$ grid whose
only interior point, $p=0.1$, is not realisable on a population of sixteen; the data
could not support the reported precision, and the data-quality audit described in
Part~A found that three of the six re-run conditions were additionally corrupted by
silent API failure. We re-ran all five lineages clean at $N=64$ with seven resolving
points inside the transition, which is the basis for item~8; the apparent spread of
$0.084$--$0.118$ was resolution, and at $N=64$ the lineages agree. \emph{(ii)} Our
capability probe covered 2--4B and 70B-plus, and the ${\sim}30$B figure interpolated
across the range between them. We have now measured that range: across fourteen
models from 2B to 123B the floor sits between 3B and 8B, a 4B-active
mixture-of-experts model passes where dense 3B models fail, and a 123B model scores
below an 8B one. The structural claim---that fitness to govern is a threshold rather
than a smooth scale curve---is unchanged and now rests on fourteen models rather than
on an interpolation; what changes is its operational form, since an admission rule
stated as a parameter count errs in both directions. \emph{(iii)} Our verifier was
free to emit any answer, including one no panel member had proposed, which makes
verifier quality unidentifiable: the same judge scored equivalently with no panel at
all. Item~5 supplies the constrained instrument and the control that separates
selection from generation, and the verifier figures are reported under it.

In each case the theory the paper advances is unchanged; the measurements behind it
are better.

\section*{Part C --- Presentation}

Following your suggestion, the body carries a compact methods overview and the
detailed protocols sit in Appendix~B. To keep the body readable at 20 pages we have
also moved to a supplementary appendix the material that supports but is not central
to the argument: the mechanical study of institutional decision rules, the full
institutional-ladder arms, the granularity account of commons survival, the
sustainability--equity trade-off, the token-budget caution, and five single-result
figures whose values are stated in the text. Each has a pointer from the body, and
no result has been removed.

We have also corrected a figure that had become inconsistent with the revised text:
the verifier figure was generated from the original free-form-judge run and depicted
the very artifact the section now disproves. All figures are now generated from the
released result files by a script in the repository, and a companion script checks
every number in the manuscript against those files and flags any figure older than
the data it depicts. The reported total API expenditure has been updated from
${\sim}\$16$ to ${\sim}\$82$, summed from the per-run cost fields in the released
results.

\section*{Summary of changes}

\begin{longtable}{@{}p{7.4cm}p{4.2cm}p{3.4cm}@{}}
\toprule
\textbf{Change} & \textbf{Section} & \textbf{Nature} \\
\midrule
\endhead
Methods section added; detailed protocols in appendix & \emph{Methods}; \emph{Appendix B} & As requested \\
Prompt templates, parameter table & \emph{Appendix B} & As requested \\
Data-quality auditing described and applied to new runs & \emph{Methods}; \emph{Threats to Validity} & New control \\
Capture threshold re-run at $N=64$ (7 resolving points) & \emph{Capture Threshold Transfers} & Central result strengthened \\
Capability floor shown to be role-specific; re-measured across 14 models & \emph{Capability Gate is a Floor} & New result; supersedes ${\sim}30$B \\
Constrained-selection instrument and the 0/112 control & \emph{Only a Constrained Verifier} & New control; supersedes earlier ceiling figure \\
$p_c(H)$ horizon law derived & \emph{Cooperation} & New law; recovers our ${\sim}25\%$ \\
Bistability analysis; content-independence experiment & \emph{Cooperation}; \emph{Threshold Does Generalize} & New result; confirms the generalization \\
Same-model head-to-head design & \emph{Cost vs.\ Quality} & Stricter test; $2.53\times$ confirmed \\
Coverage precondition and competence matching & \emph{Coverage Precondition} & New result \\
Verification-easier-than-generation ladder & \emph{Only a Constrained Verifier} & New result \\
Institutional ladder & \emph{Cooperation} & New result \\
Ostrom (1990) added & References & New citation \\
Figures regenerated from result files; two were stale & --- & Consistency fix \\
Peripheral material moved to supplement & \emph{Appendix C} & Presentation \\
\bottomrule
\end{longtable}

A version with all changes in coloured text accompanies this submission, as
requested in the decision letter.

We recognise this is more new work than a Methods request would normally produce.
Having had to specify our procedures precisely, it seemed better to test the claims
that specification put in question than to leave them resting on the resolution we
originally had. The paper's theoretical contribution---the competence band, the
capture threshold, the capability floor---is unchanged; what has changed is the
quality of the evidence behind it, and the addition of four findings we did not have
before. We are grateful for a review process that improved the work rather than
merely judging it, and we hope the revised manuscript is now in a form your referees
can assess.

\bigskip
Sincerely,

\medskip
Gopal Kalyanaraman and Vijay K.\ Madisetti

\end{document}
